\documentclass[a4paper, oneside, 12pt]{article}

% Paquets pour le Français
\usepackage[utf8]{inputenc} % Gestion encodages
\usepackage[T1]{fontenc} % ???
\usepackage[francais]{babel} % Typographie française

% Images
\usepackage{graphicx}

% Mise en page
\usepackage[top=2cm, bottom=2cm, left=2cm, right=2cm]{geometry}
\usepackage{appendix}
% Outils pour l’écriture
\usepackage{xcolor}
\usepackage{layout}
\usepackage{lipsum}

\newcommand{\todo}[1][Il y a là encore des choses à écrire !]{\colorbox{yellow}{#1}\\}

\newcommand{\sectionSpeciale}[1]{\section*{#1}\addcontentsline{toc}{section}{#1}}
\newcommand{\sectionCachee}[1]{\addcontentsline{toc}{section}{#1}}

%%%%%%%%%%%%%%%%%%%% Gestion des annexes %%%%%%%%%%%%%%%%%%%%
\makeatletter

	\newcommand{\annexe}[1]{\addcontentsline{ann}{section}{#1}\newpage\section*{#1}}
	\newcommand\toann{\sectionSpeciale{Annexes}\input{annexes.tex}}
      \newwrite\tf@ann
      \immediate\openout\tf@ann\jobname.ann\relax
      
\makeatother

%%%%%%%%%%%%%%%%%%%% Abréviations %%%%%%%%%%%%%%%%%%%%

% TODO
\newcommand\bash[1]{\texttt{#1}}

\newcommand\kam{Kamailio}
\newcommand\frad{FreeRADIUS}

%%%%%%%%%%%%%%%%%%%% Début du document %%%%%%%%%%%%%%%%%%%%
\begin{document}

\renewcommand{\contentsname}{Table des matières}
\renewcommand{\listfigurename}{Table des illustrations}
\renewcommand{\figurename}{{\sc Illustration}}

%%%%%%%%%%%%%%%%%%%% Page de garde %%%%%%%%%%%%%%%%%%%%

	% Pas de numéro sur cette page
	\thispagestyle{empty}

	% Note : logo disponibles sur http://fc.isima.fr/~brunot/files/isima/logos/
	% J’avais trouvé une version vectorielle, ce serait mieux, mais je ne la trouve plus…
	\includegraphics[width=6cm]{isima.png}
	
	\vspace{0.5cm}
	
	\begin{minipage}{4cm}
	\begin{flushleft}
		Institut Supérieur d’informatique, de Modélisation et de leurs Applications
		
		\vspace{0.5cm}
		
		\small{ Campus des Cézeaux \\ 24 rue des Landais \\ BP 10125 \\ 63173 Aubière CEDEX }
	\end{flushleft}
	\end{minipage}
	
	\vspace{4cm}

	\begin{center}
		Rapport d’ingénieur \\
		Projet de 2{\ieme} année \\
		Filière {\em{Réseaux et télécommunications}} \\
		\Large{Facturation d’appels sur un serveur Kamailio/Radius}
	\end{center}
	
	\vspace{4cm}
	
	% TODO %

\newpage
	
%%%%%%%%%%%%%%%%%%%% Précontenu %%%%%%%%%%%%%%%%%%%%

\pagenumbering{roman}

%%%%%%%%%% Remerciements %%%%%%%%%%

\sectionSpeciale{Remerciements}

\newpage

%%%%%%%%%% Table des illustrations %%%%%%%%%%

\sectionCachee{\listfigurename}

\listoffigures

\newpage

%%%%%%%%%% Résumé/Abstract %%%%%%%%%%

\sectionSpeciale{Résumé}

\vspace{10cm}

\sectionSpeciale{Abstract}

\newpage

%%%%%%%%%% Table des matiètes %%%%%%%%%%

\sectionCachee{\contentsname}

\tableofcontents

\newpage

%%%%%%%%%%%%%%%%%%%% Contenu %%%%%%%%%%%%%%%%%%%%

% Sauvegarde du numéro de la page courante
\newcounter{metapage}
\setcounter{metapage}{\value{page}}
% Remise à zéro des numéros de page & changement de chiffres
\pagenumbering{arabic}

%%%%%%%%%% Introduction %%%%%%%%%%

\sectionSpeciale{Introduction}

Ceci est une introduction

\newpage

% On verra le placement plus tard…
\section{Installation}

\subsection{Installation du système}

\todo[Trouver la version de Debian]

La première étape de notre projet consistait en l’installation de serveurs de téléphonie SIP. Nous en avons installé deux, fontionnant avec le système d’exploitation Debian. Conformément à notre cahier des charges, nous y avons installé un proxy SIP, {\kam}, attaché à un serveur d’authentification {\rad}, {\frad}.

\subsection{\kam}

\subsubsection{Installation}

Pour installer {\kam}, il nous a fallu commencer par installer les dépendances qui n’étaient pas encore présentes dans le système, disponibles dans les paquets :

\begin{itemize}
	\item{gcc} ;
	\item{bison} ;
	\item{libmysqlclient-dev} ;
	\item{libradiusclient-ng-dev} ;
	\item{mysql-server} ;
	\item{mysql-client} ;
	\item{make}.
\end{itemize}

Nous avions initialement envisagé d’utiliser ensuite les paquets Debian disponibles sur les serveurs de {\kam}, mais le paquet de base n’est pas compilé avec le support de {\rad}, rendant le paquet des bibliothèques d’authentification {\rad} inopérant.

Nous avons donc téléchargé et décompressé les sources de {\kam} pour les compiler. Nous avons choisi d’utiliser la dernière version disponible, qui était la version 4.2.2. Pour cela, nous avons utilisé les commandes :

\begin{shellcode}
wget http://www.kamailio.org/pub/kamailio/4.2.2/src/kamailio-4.2.2_src.tar.gz
tar xvf kamailio-4.2.2_src.tar.gz
\end{shellcode}

Puis nous sommes allés dans le dossier ainsi créé, pour notre part kamailio-4.2.2 et nous avons exécuté :

\begin{shellcode}
make cfg
\end{shellcode}

Cette commande construit la configuration par défaut pour la compilation, qu’il a fallu adapter à nos besoins. Conformément aux indications du rapport de l’année dernière, nous avons effectué les modifications suivantes :

\begin{itemize}
	\item{dans la liste de modules exclus du fichier \texttt{module.lst}, nous avons retiré les noms des modules concernant {\my} et {\rad}}.
\end{itemize}

Nuos avons ensuite compilé et installé {\kam} :
	
\begin{shellcode}
make all
make install
\end{shellcode}

Cette dernière commande installe {\kam} dans \texttt{/usr/local}.

\subsubsection{Configuration}

Nous avons dû ensuite configurer {\kam}. Nous sommes allés dans le dossier \texttt{/usr/local/etc/kamailio}. Il a fallu rajouter les lignes suivantes dans le fichier de configuration \texttt{kamailio.cfg} :

\#!define WITH\_MYSQL

\#!define WITH\_AUTH

\#!define WITH\_USRLOCDB

\#!define WITH\_NAT

Nous avons alors créé puis configuré la base de données. Pour la créer, il a fallu d'abord préciser quel système de base de données nous allions utiliser.
Dans notre cas, c'est mysql. C'est pourquoi il a fallu aller dans le fichier \texttt{/usr/local/etc/kamailio/kamctlrc} afin de décommenter la ligne :

DBENGINE=MYSQL

Ensuite, nous avons tapé 
\begin{shellcode}
/usr/local/sbin/kamdbctl create
\end{shellcode}
\todo

\subsection{\frad}
\subsubsection{Installation}

Nous avons choisi ensuite d'installer {\frad} à partir des paquets présents sur la machine. Pour cela, il a fallu taper la commande :

\begin{shellcode}
apt-get install freeradius
\end{shellcode}

Il fallait aussi installer les outils nécessaires au bon fonctionnement de {\frad} :

\begin{shellcode}
apt-get install freeradius-utils
\end{shellcode}

Puis, pour finir, nous avons installé un client Radius nommé Radiusclient grâce à la commande suivante :

\begin{shellcode}
apt-get install radiusclient1
\end{shellcode}

\subsubsection{Configuration}
Nous avions fini d'installer {\frad}. Il nous restait donc à le configurer. Ce serveur {\rad} peut être configuré de deux manières différentes. Soit à partir des fichiers présents dans \texttt{/etc/freeradius}, soit à partir d'une base de données. Nous avons choisi de configurer à partir des fichiers. Il a donc fallu aller vérifier si la configuration se faisait bien grâce aux fichiers. Pour cela, il a suffi d'aller vérifier dans le fichier \texttt{/etc/freeradius/sites-available/default}.

Ensuite, nous sommes allés dans le répertoire \texttt{/etc/freeradius} afin de modifier le fichier \texttt{users}. Dans ce fichier, nous avons ajouté des utilisateurs de la façon suivantes :

login Cleartext-Password:="password"

Cette ligne autorise l'utilisateur login à se connecter avec le mot de passe "password".

Nous avons pû ensuite faire le test en utilisant la commande suivante :

\begin{shellcode}
radtest login password localhost 0 testing123
\end{shellcode}

Le paramètre "testing123" est le mot de passe donné par défaut au client localhost dans le fichier \texttt{/etc/freeradius/client.conf}. Nous l'avons ensuite changé en "kamailio".


\subsubsection{Intégration de {\rad} à {\kam}}

Pour commencer, il a fallu copier les dictionnaires \texttt{/usr/local/etc/kamailio/dictionary.kamailio} et \texttt{/usr/local/etc/kamailio/dictionary.sip} dans le répertoire \texttt{/etc/radiusclient-ng}.

\begin{shellcode}
cp /usr/local/etc/kamailio/dictionary.kamailio /etc/radiusclient-ng

cp /usr/local/etc/kamailio/dictionary.sip /etc/radiusclient-ng
\end{shellcode}

Puis, nous avons inclus dans \texttt{/etc/radiusclient-ng/dictionary} les deux fichiers suivants :

\begin{itemize}
\item{\texttt{/etc/radiusclient-ng/dictionary.kamailio}}
\item{\texttt{/etc/radiusclient-ng/dictionary.sip}}
\end{itemize}

Puis il a fallu décommenter les lignes des attributs ou des valeurs Failed, Sip-Session et Call-Check dans les trois fichiers dictionary.

Nous avons ensuite dû vérifier que les deux lignes suivantes,qui donnent l'adresse du serveur {\rad}, se trouvaient bien dans le fichier \texttt{/etc/radiusclient-ng/radiusclient.conf} :

\begin{itemize}
\item{authserver localhost}
\item{acctserver localhost}
\end{itemize}

Nous avons ensuite été modifier le fichier \texttt{/etc/radiusclient-ng/servers}. Nous avons annoncé notre serveur de la façon suivante :


192.168.102.120		kamailio


Ce fichier contient l'adresse des serveurs {\rad} et leurs clefs. Ainsi nous avons choisi la clef "kamailio".

Puis ensuite il a fallu aller modifier le fichier \texttt{/etc/freeradius/clients.conf} qui définit les clients autorisés par le serveur {\frad} :

\begin{shellcode}
Client localhost
{
 ipaddress = 192.168.102.120
 secret = kamailio
}
\end{shellcode}

Nous avons dû ensuite modifier le fichier \texttt{/etc/freeradius/users}. Nous avons tout effacer et avons écrit les utilisateurs autorisés ou pas de la manière suivante :


7000@192.168.102.120 Auth-Type=Accept
7001@192.168.102.120 Auth-Type=Accept
7003@192.168.102.120 Auth-Type=Reject


Cela autorise les utilisateurs 7000 et 7001 mais rejette l'utilisateur 7002.

Ainsi nous avons deux sécurités à passer afin de pouvoir utiliser le serveur {\kam}. Il faut tout d'abord être dans la base de donnée de {\kam} puis ensuite dans le fichier users de {\frad} en tant qu'utilisateur accepté.




\newpage

%%%%%%%%%% Conclusion %%%%%%%%%%

\sectionSpeciale{Conclusion}

\newpage

%%%%%%%%%%%%%%%%%%%% Postcontenu %%%%%%%%%%%%%%%%%%%%

% Changement des chiffres utilisés pour les numéros de pages
\pagenumbering{roman}
% Rétablissement du numéro de page méta
\setcounter{page}{\value{metapage}}

%%%%%%%%%% Références bibliographiques %%%%%%%%%%

\sectionSpeciale{Références bibliographiques}

\newpage

%%%%%%%%%% Annexes %%%%%%%%%%

\toann

\annexe{Une annexe}
\annexe{Une autre annexe}
\annexe{Encore une annexe}

%%%%%%%%%%%%%%%%%%%% Fin %%%%%%%%%%%%%%%%%%%%

\end{document}
